\documentclass[conference]{IEEEtran}
\IEEEoverridecommandlockouts
\usepackage{cite}
\usepackage{amsmath,amssymb,amsfonts}
\usepackage{graphicx}
\usepackage{textcomp}
\usepackage{xcolor}
\usepackage[english]{babel}
\usepackage{listings}
\usepackage{booktabs}
\usepackage{hyperref}

\lstset{
  basicstyle=\ttfamily\footnotesize,
  breaklines=true,
  frame=single,
  keywordstyle=\bfseries,
  commentstyle=\color{gray}\itshape,
  showstringspaces=false,
  tabsize=4,
  morekeywords={func,type,struct,chan,go,select,case,return,if,for,range,
  make,nil,true,false,var,const,package,import,defer,error},
}

\def\BibTeX{{\rm B\kern-.05em{\sc i\kern-.025em b}\kern-.08em
T\kern-.1667em\lower.7ex\hbox{E}\kern-.125emX}}

\begin{document}

\title{Distributed Configuration Updater:\\
  Design and Implementation of a CRDT-Based\\
Settings Replication System}

\author{\IEEEauthorblockN{Gabriele Da Re}
  \IEEEauthorblockA{\textit{Dept. Mathematics} \\
    \textit{Universit\`{a} degli Studi di Padova}\\
    Padova, Italy \\
  gabriele.dare00@gmail.com}
}

\maketitle

\begin{abstract}
  This paper describes the architecture, implementation and decisions made during
  the creation of a configuration synchronization system for server clusters.
  The system synchronises Key-Value (KV) entries across nodes through: Last-Writer-Wins (LWW),
  Conflict-free Replicated Data Type (CRDT), gRPC transport protocol and a Hybrid Logical Clock (HLC) for causal ordering.
  We explain the {\color{violet}modules design decisions} and the problems encountered/mitigated during development such as: Clock-skew,
  crash persistence handling, tombstones management and the anti-entropy approach and lifecycle enforcement between initialisation
  and runtime phases.
\end{abstract}

\begin{IEEEkeywords}
  CRDT, distributed systems, data consistency, hybrid logical clocks,
  gRPC, Go, key-value store
\end{IEEEkeywords}

% ─────────────────────────────────────────────────────────────────────────────
\section{Introduction}
% ─────────────────────────────────────────────────────────────────────────────

Configuration management across a server cluster is a common operational
challenge. Changes to settings such as feature flags, service endpoints,
parameters must propagate to every node without a coordinator becoming
the single failure point.

Approaches such as \href{https://etcd.io/}{etcd},
\href{https://developer.hashicorp.com/consul}{Consul} rely on a leader node or a
distributed database. These provide linearisability at the cost of availability
({\color{violet}Consistency and Partition tolerance part of the CAP theorem\cite{gilbert2012}}; ACID principle): since RAFT\cite{raft}
based, when the leader or the majority ($N/2+1$ nodes) is unreachable,
there is no working system (not always Available; CAP\cite{gilbert2012}). For configuration data, human operators can intervene
on conflicts. Moreover, having stale data in some circumstances is relatively low risk.
{\color{violet}Therfore}, the trade-off favours Availability and Partition tolerance {\color{violet}
(AP\cite{gilbert2012})} {\color{violet}at the cost of strong} Consistency.

Conflict-free Replicated Data Types (CRDTs) \cite{shapiro2011} offer an
alternative: every node can accept writes locally, the {\color{violet}CRDT, a join-semilattice under merge,} guarantees that any two nodes with the same set of updates will
converge to the same state without inter-node coordination. This paper describes a
practical implementation of a CRDT based configuration store targeting small,
fixed-membership clusters where $O(100)$ configuration keys are replicated
across $O(10)$ nodes. This limit can be pushed up to $O(1000)$ keys and $O(50)$ nodes
if necessary, based upon limits imposed by the full snapshot sync implementation (Section~\ref{sec:antientropy}).

The implementation is written in Go and uses gRPC for inter-node
communication. The remainder of this paper is structured as follows.
Section~\ref{sec:arch} presents the overall system architecture.
Section~\ref{sec:crdt} describes the CRDT data model.
Section~\ref{sec:hlc} covers the hybrid logical clock.
Section~\ref{sec:actor} explains the actor-based concurrency model.
Section~\ref{sec:network} details the networking module.
Section~\ref{sec:antientropy} discusses the periodic state reconciliation.
Section~\ref{sec:persistence} describes crash-safe persistence.
Section~\ref{sec:gc} presents tombstone garbage collection.
Section~\ref{sec:lifecycle} covers lifecycle and error handling.
Section~\ref{sec:testing} outlines the testing strategy.
Section~\ref{sec:conclusion} presents the conclusions.

% ─────────────────────────────────────────────────────────────────────────────
\section{System Architecture}
\label{sec:arch}
% ─────────────────────────────────────────────────────────────────────────────

The code is structured as five Go packages with a strict dependency
hierarchy, shown in Table~\ref{tab:packages}.

\begin{table}[htbp]
  \caption{system module breakdown}
  \label{tab:packages}
  \begin{center}
    \begin{tabular}{ll}
      \toprule
      \textbf{Package} & \textbf{Responsibility} \\
      \midrule
      \texttt{types}      & Shared domain types: HLC, SettingEntry, Snapshot \\
      \texttt{crdt}       & LWW CRDT, state machine, persistence \\
      \texttt{logic}      & Local settings file watcher and writer \\
      \texttt{network}    & gRPC server, streaming client, proto conversion \\
      \texttt{controller} & Wires all packages, lifecycle and backgrounds tasks \\
      \midrule
      \texttt{main}       & CLI entry-point \\
      \bottomrule
    \end{tabular}
  \end{center}
\end{table}

The data flows through the system following four paths, wired up in
\texttt{controller.Run}:

\begin{enumerate}
  \item \textbf{Local write}: the \texttt{logic} watcher detects file
    changes on write and pushes a \texttt{SettingEntry} to \texttt{crdt.NotifyLocal},
    which queues it in local channel of the CRDT actor. The CRDT ticks
    its clock, timestamps the entry and merges it, persists state, and calls
    \texttt{OnBroadcast} delegating the broadcast to the network layer.
  \item \textbf{Remote write}: a peer to which the node is subscribed streams an update,
    \texttt{network.Client} calls \texttt{crdt.NotifyRemote}. The CRDT
    updates its clock, merges, persists, and calls \texttt{OnFileSync}
    that delegates the \texttt{logic} layer to write the change locally.
  \item \textbf{Anti-entropy}: a ticker in the controller calls \texttt{syncAllPeers},
    which sends the local snapshot to each peer via a Sync RPC which returns (if any)
    any newer entries through \texttt{NotifyRemote}.
  \item \textbf{Tombstone GC}: another separate ticker calls
    \texttt{crdt.PurgeTombstones}, which removes stale deleted entries
    from the in-memory state map via the CRDT actor model.
\end{enumerate}

The controller is the only package that depends on all the others, preventing import cycles.

% ─────────────────────────────────────────────────────────────────────────────
\section{CRDT Data Model}
\label{sec:crdt}
% ─────────────────────────────────────────────────────────────────────────────

\subsection{Data Model}

{\color{violet}The data structure is an \textbf{LWW-register}\cite{shapiro2011} map keyed by setting name.
  Per Shapiro, each register is any local data type $X$ paired with its timestamp from the $assign$ function.
  In our case $X$ is a $(key, value, Deleted)$ triplet, and the timestamp is managed through an HLC (\autoref{sec:hlc}).
  The key is stored alongside its value so the map can be rebuilt on the destination node, whilst the \texttt{Deleted} flag deals with
  deleted entries on the setting file, enabling the local module to drop locally persisted settings.
}

\begin{lstlisting}
type SettingEntry struct {
    Key     string
    Value   string
    Clock   HLC
    Deleted bool
}
\end{lstlisting}

\subsection{Merge Semantics}
\label{merge-semantics}

Two nodes merge by iterating {\color{violet}the local map against the remote one and checking each key:
if the key exists locally with a higher clock is kept unaltered, else it is created or overwritten.}
On equal clocks, the lexicographically larger NodeID is used as tiebreaker,
ensuring a deterministic total order.

\begin{lstlisting}
func (c *CRDT) merge(incoming SettingEntry) bool {
    existing, ok := c.state[incoming.Key]
    if ok && !existing.Clock.Before(incoming.Clock) {
        return false
    }
    c.state[incoming.Key] = incoming
    return true
}
\end{lstlisting}

\subsection{Tombstone Semantics}

Deletions are {\color{violet}managed by setting the \texttt{Deleted} flag to \texttt{true}
and stamping a fresh clock on the associated triplet}. Tombstones undergo the same LWW merge as the payload and
can be displaced by a later write. This avoids the anomaly where a deletion
arrives after a concurrent write and erroneously removes the newer value.

\subsection{CvRDT}

The implementation is {\color{violet}fundamentally a state-based CRDT: a timed
  \texttt{Sync} provides anti-entropy reconciliation and each local state change broadcasts a
  \texttt{SettingEntry} to peers (operation-based style).
  Because LWW merge is commutative and associative, a set of updates can be applied regardless of the order to obtain the same result,
and idempotent, makes a merge of $B\subseteq A$ a no-op. Therfore,} each broadcast is algebraically
equivalent to a partial state merge. Convergence is guaranteed
entirely by the state-based anti-entropy \texttt{sync} cycle, the broadcast is a best-effort optimisation.

\subsection{CvRDT proof}
\label{sec:cvrdtproof}

\paragraph{State forms a join-semilattice.}
Define the replica state as $S = \mathrm{Map}[\mathit{String} \to
\mathit{SettingEntry}]$.

\noindent Impose a partial order $\forall k \in Keys$
\begin{equation}
  S_1 \leq S_2 \iff \forall k: S_1[k].\mathit{Clock} \leq S_2[k].\mathit{Clock}
\end{equation}
where \texttt{Before} defines the strict order $<$ on clocks (Definition~\ref{sec:hlc}),
and $\leq$ is its reflexive closure: $a \leq b \iff a < b \;\lor\; a = b$.
The least upper bound per key is
\begin{equation}
  (S_1 \sqcup S_2)[k] = \arg\max_{\mathit{Clock}}\bigl(S_1[k],\; S_2[k]\bigr)
\end{equation}
which is exactly \texttt{merge()}.

\paragraph{Merge implements the LUB.}
Let $A$, $B$, $C$ be \texttt{SettingEntry} values for the same key.

\begin{itemize}
  \item \textbf{Commutative:} $\mathrm{merge}(A,B) = \mathrm{merge}(B,A)$, since
    $\max(\mathit{Clock}_A, \mathit{Clock}_B)$ is symmetric.
  \item \textbf{Idempotent:} $\mathrm{merge}(A,A) = A$, since
    $\max(T,T) = T$.
  \item \textbf{Associative:} $\mathrm{merge}(A,\mathrm{merge}(B,C)) = \mathrm{merge}(\mathrm{merge}(A,B),C)$, since $\max$ is associative.
\end{itemize}

\paragraph{Monotonicity.}
Every merge moves state upward in the lattice:
\begin{equation}
  S \leq \mathrm{merge}(S,\, e),
\end{equation}

since \texttt{merge} replaces an entry only when the incoming clock is strictly
greater, and never decreases a clock value.

\paragraph{Liveness.}
\texttt{syncAllPeers} runs periodically exchanging full snapshots. After a
network partition, reconnection triggers \texttt{Sync}, delivering all missing updates.
Hence every update eventually reaches every replica.

\paragraph{Conclusion.}
By the Eventual Consistency theorem for CvRDTs \cite{shapiro2011}: any
two replicas that have received the same set of updates hold identical maps,
regardless of delivery order.

% ─────────────────────────────────────────────────────────────────────────────
\section{Hybrid Logical Clock}
\label{sec:hlc}
% ─────────────────────────────────────────────────────────────────────────────

\subsection{Motivation}

A Lamport clock \cite{lamport1978} provides causal ordering but loses the physical-time
locality property: events that happen close together in real time may receive
widely separated clock values if logical increments accumulate. Hybrid Logical
Clocks \cite{kulkarni2014} integrate the Lamport clock with a wall time,
preserving the benefits of both.

\subsection{Clock Structure}

{\color{violet}Each node maintains an HLC as a triple $(wall, logical, nodeID)$}. The wall field stores
Unix nanoseconds, the logical field is a 32-bit logical counter, the NodeID is a UUID
string used as tiebreaker (Section~\ref{merge-semantics}). This ensures \texttt{Before()}
commutativity since $\forall n_i, n_j\in Nodes,\; i\neq j: Id(n_i) \neq Id(n_j)$

\subsection{Tick and Update Rules}

On a local event (\emph{tick}), the clock advances ($\ell_l, \ell_r$ for
$logical_{local}, logical_{received}$, $w_l, w_r$ for $wall_{local}, wall_{received}$):
\begin{align}
  wall' &= \max(w_l,\; now), \notag \\
  logical' &=
  \begin{cases}
    logical + 1 & wall' = w_l \\
    0           & otherwise
  \end{cases}
\end{align}

On receiving a remote clock (\emph{update}), the canonical
Kulkarni--Demirbas rule \cite{kulkarni2014} is applied:
\begin{equation}
  wall' = \max(w_l,\; w_r,\; now)
\end{equation}
The logical counter then captures concurrency at the chosen wall time:
\begin{equation}
  logical' =
  \begin{cases}
    \max(\ell_l, \ell_r) + 1 & w' = w_l = w_r \\
    \ell_r + 1                 & w' = w_r \neq w_l \\
    0                          & otherwise
  \end{cases}
\end{equation}

\subsection{Clock Skew Handling}

The system assumes NTP-disciplined clocks with bounded skew ($\lesssim$10ms on
  LAN, $\lesssim$100ms on WAN. CockroachDB's empirically chosen tolerance is
500ms \cite{cockroachdb2023}). Under this assumption HLC self-corrects: when a
node physical clock falls behind the cluster wall time, the logical counter
increments, preserving causality without advancing the wall.

An initial design rejected incoming entries whose \texttt{WallTime} exceeded the
receiver's physical clock by more than a fixed threshold $T$:

\begin{equation}
  \textit{received.WallTime} - \textit{now} > T \implies \text{reject}
\end{equation}

This check merges the HLC wall with the sender's physical clock.
Consider a cluster at physical time $0$, a fast
node $B$ at $+0.9s$, and a slow node $A$ at $-0.9s$
(so $|t_A - t_B| = 1.8s> T$):

\begin{enumerate}
  \item Cluster accepts $B$: $\delta = 0.9s < T$.  HLC wall bumped to $+0.9s$.
  \item Cluster sends its own entries to $A$.  Those entries carry
    $\texttt{WallTime} = +0.9s$.
  \item $A$ computes $0.9 - (-0.9) = 1.8s > T$ and rejects,
    despite $A$ and the cluster being only $0.9s$ apart physically.
\end{enumerate}

This produces a \textbf{self-healing partition}: $A$ stops
receiving updates despite full network connectivity. The partition resolves
automatically once NTP corrects $A$'s clock past the acceptance threshold
($\textit{now}_A > \textit{HLC}_\text{cluster} - T$), after which the next
anti-entropy cycle delivers all missed entries and convergence resumes. However,
NTP slews clocks gradually (does not step unless the offset exceeds $\approx
128ms$ by default), so the partition may persist for minutes to hours. If both
$A$ and the cluster are stuck with bad NTP indefinitely the partition does not
heal at all.

The CRDT merge function remains correct so safety is preserved; the failure point
is at the delivery layer, violating the liveness precondition of Strong
Eventual Consistency \cite{shapiro2011}.

Note that the check could have been maintained if: (i) the check was performed on
the sender's physical clock rather than the HLC wall, so each node judges independently;
(ii) the threshold was set to a greater value such as 30s, to detect
real skew problems while avoiding triggers on NTP jitter.
In the current implementation, skew check is removed, protection is
delegated entirely to NTP and operational monitoring.

% ─────────────────────────────────────────────────────────────────────────────
\section{Actor Model}
\label{sec:actor}
% ─────────────────────────────────────────────────────────────────────────────

All mutations to the CRDT state map are serialised through a single actor
goroutine running a \texttt{select} loop over five typed channels:

\begin{lstlisting}
type CRDT struct {
    state      map[string]SettingEntry
    localCh    chan SettingEntry
    remoteCh   chan SettingEntry
    queryCh    chan queryRequest
    snapshotCh chan snapshotRequest
    gcCh       chan int64
}
\end{lstlisting}

\textbf{No mutex} is required: only the actor goroutine accesses channels and
CRDT state ensuring operation atomicity. External callers (\texttt{NotifyLocal}, \texttt{NotifyRemote},
\texttt{Snapshot}, \texttt{queryRequest}, \texttt{PurgeTombstones}) communicate exclusively through
channels. Snapshot and query requests are buffered synchronous response channels thus blocking
until the request is serviced.

Callbacks \texttt{OnBroadcast}, \texttt{OnFileSync} are invoked
synchronously in the actor loop {\color{violet}to preserve event order. \texttt{OnFileSync}
  decouoples the actor loop from disk I/O, queuing each entry on a buffered channel drained by
calling in order the \texttt{Write} function for each entry by a writer goroutine.}

% ─────────────────────────────────────────────────────────────────────────────
\section{Networking Module}
\label{sec:network}
% ─────────────────────────────────────────────────────────────────────────────

\subsection{gRPC Service}

The inter-node protocol exposes two RPCs over gRPC:

\begin{itemize}
  \item \textbf{SubscribeStateUpdates} (server-streaming): the caller receives
    \texttt{ServerStateUpdate} messages through a persistent stream every time
    the server calls \texttt{Broadcast}. This ensures real-time propagation of
    local edits.
  \item \textbf{Sync} (unary): the caller sends its local snapshot, the
    server updates its entries and responds with local server side entries that
    are newer. Used for anti-entropy reconciliation.
\end{itemize}

The server maintains a map of subscriber channels protected by a mutex.
\texttt{Broadcast} under read-lock retrieves the subscribed peers channels, then
sends the update to each channel in a non-blocking manner. Slow subscribers
drop updates relying on the next anti-entropy pass to recover.

\subsection{Abstraction Boundary}

A design decision was taken such that protobuf conversions
(\texttt{ToProto}, \texttt{FromProto}, \texttt{SnapshotToProto}) are
confined to the \texttt{network} package. The \texttt{Client} methods accept
and return domain types (\texttt{types.Snapshot}, \texttt{types.SettingEntry}).
Conversion to and from wire format is performed internally.

This was not planned in advance: an earlier iteration of \texttt{Client.Sync} and \texttt{Client.Broadcast}
accepted \texttt{[]*userpb.SettingEntry} and \texttt{[]*userpb.ServerStateUpdate},
thus requiring the controller to perform the conversion before each call.
Refactor moved conversion dependency into network package, making the transport
layer a true black box.

\subsection{Reconnection and Keepalives}

{\color{violet}A \texttt{Subscribe} method defined in \texttt{Client} wraps the streaming RPC in
  a retry loop as follows:

\begin{lstlisting}
  func (c *Client) Subscribe(
    ctx context.Context,
    getSnapshot func() types.Snapshot,
    onUpdate func(types.SettingEntry)) {
    for {
      s := getSnapshot()
      if err := c.runStream(ctx, s, onUpdate); err != nil {
        if ctx.Err() != nil {
          return
        }
        log.Printf("stream broken, reconnecting: %v", err)
        time.Sleep(3 * time.Second)
      }
    }
}
\end{lstlisting}

  On any error (unless cancelled context), including server restart or network
  interruption, it waits three seconds and tries to re-establish the stream through \texttt{runStream}:
  it first performs a Sync call to replay any entries missed during the outage, then opens the update stream.
gRPC keepalive parameters are configured to detect dead connections.}

% ─────────────────────────────────────────────────────────────────────────────
\section{Anti-Entropy Reconciliation}
\label{sec:antientropy}
% ─────────────────────────────────────────────────────────────────────────────

Broadcast over gRPC is best-effort: a slow, offline or reconnecting subscriber
may miss updates. A (state-based) CRDT that relies solely on broadcast for
propagation is not eventually consistent in presence of message losses,
nodes that miss a broadcast permanently diverge.

Production CRDT deployments (Riak, Dynamo \cite{dynamo2007}) address convergence
with a background anti-entropy mechanism: periodic full-state comparison
identifies and repairs divergence independently of the broadcast path (\href{https://docs.riak.com/riak/kv/latest/learn/concepts/active-anti-entropy/index.html}{Riak AAE}).

Our system implements a lightweight variant: \texttt{syncAllPeers}
is called periodically by the controller; this sends a local snapshot to
each peer via the Sync RPC. The peer merges any missing or newer entries and returns the entries where its copy is newer
or not present on caller.
This is sufficient for eventual consistency: even if every broadcast drops, {\color{violet}the
state will eventually converge.}

There is an unimplemented enhancement regarding exchanges of redundant KV pairs: it
can happen that two nodes send a Sync request {\color{violet}concurrently}; each node takes a snapshot
of local state and sends it, then upon receiving the snapshot both nodes update their local entries and send
the updates they consider necessary to the other, which has already applied them through the received snapshot.
This is solvable but since it is a performance issue and not a correctness one, it is not covered in this
project (e.g. Jitter, follower/leader Sync, Merkle tree, Delta state, etc.).

\begin{lstlisting}
func (ctrl *Controller) syncAllPeers(ctx context.Context) {
    snap := ctrl.crdt.Snapshot()
    for _, client := range ctrl.clients {
        go client.Sync(ctx, snap, func(e SettingEntry) {
            ctrl.crdt.NotifyRemote(e)
        })
    }
}
\end{lstlisting}

A 60-second interval was chosen as default for the sync anti-entropy cycle.

% ─────────────────────────────────────────────────────────────────────────────
\section{Crash-Safe Persistence}
\label{sec:persistence}
% ─────────────────────────────────────────────────────────────────────────────

After each merge, CRDT serialises its state to \texttt{crdt\_state.json} found in the CRDT
working directory. The naive approach using \texttt{os.WriteFile} is not crash-safe:
if the process is killed after \texttt{WriteFile} had truncated the file and before
the kernel had flushed all pages to disk, the node will start with a corrupted state file.

The correct pattern we used is the atomic rename:

\begin{lstlisting}
func atomicWriteFile(path string, data []byte) error {
    tmp, err := os.CreateTemp(filepath.Dir(path), ".tmp_*")
    // ... write data to tmp, call tmp.Sync(), tmp.Close() ...
    return os.Rename(tmp.Name(), path)
}
\end{lstlisting}

Three properties make \texttt{atomicWriteFile} safe:
\begin{enumerate}
  \item The temporary file is in the same directory of the target,
    thus the rename is within a single filesystem so, atomic
    at the Virtual File System layer on POSIX systems.
  \item \texttt{tmp.Sync()} flushes file data to storage before
    the rename, preventing data from remaining only in the page cache during power-loss.
  \item \texttt{os.Rename} is atomic on POSIX: any reader will see either the
    old file or the new file, never a partially-written intermediate state.
\end{enumerate}

If the process crashes before the rename, the temporary file is
left in the directory and cleaned up during the next \texttt{crdt.Init()} call,
the original file remains intact.

On any restart, \texttt{crdt.Reconcile} is called once from \texttt{controller.Run()},
fully restoring local file integrity. This is a limiting factor since on an already initialized
node, any edits made on the settings file while the node is offline will be discarded in favour
of the last known state.

% ─────────────────────────────────────────────────────────────────────────────
\section{Garbage Collection}
\label{sec:gc}
% ─────────────────────────────────────────────────────────────────────────────

\subsection{Problem}

\textbf{Deleted entries} persist in CRDT state file indefinitely; this also
happens for \textbf{temporary files} generated by \texttt{atomicWriteFile} upon system
crash. In systems with frequent delete-recreate cycles over large KV stores,
tombstone accumulation brings two problems: memory leak over time,
slower \texttt{Sync} and \texttt{crdt.Run} since larger snapshots are required.

\subsection{Init cleanup}

With regard to orphaned tmp files from \texttt{atomicWriteFile}, the best solution in this case
was also the easiest: removal of all pattern-matched tmp files on the \texttt{crdt.Init} call.

\subsection{TTL-Based Approach}

The system also implements a TTL-based garbage collector. A ticker in the
controller calls \texttt{crdt.PurgeTombstones(ttl)} passing the TTL.
The CRDT actor then processes the request from \texttt{gcCh} channel monitored
in its Run loop, thus maintaining the single-writer on \texttt{c.state}, removing
tombstones with wallTime older than \texttt{now - ttl}.

\begin{lstlisting}
func (c *CRDT) purgeTombstones(ttl int64) {
    diff := time.Now().UnixNano() - ttl
    purged := 0
    for key, entry := range c.state {
        if entry.Deleted && entry.Clock.WallTime < diff {
            delete(c.state, key)
            purged++
        }
    }
    if purged > 0 {
        c.saveState()
        log.Printf("tombstone GC: purged %d entries", purged)
    }
}
\end{lstlisting}

The default TTL is 14 days, twice the intended maximum node downtime
after which an operator would investigate. This ensures that tombstones have been
observed and applied by all healthy nodes via anti-entropy before they are removed.
After a week of downtime a node typically requires re-initialization.

\subsection{Safety Considerations}

The TTL approach has one known failure point: if a node is offline for longer
than the TTL and holds a non-deleted entry for a key tombstoned and GC'd on all other nodes,
the node will republish the entry, so the deletion will appear to be reversed.
This is solved by forcing the node to re-initialize with empty settings if,
on startup, \texttt{time.Now() - lastShutdownTime > TTL}.
Reinit through \texttt{fixNodeState} discards any unsynced local CRDT state but cluster state is fully restored via anti-entropy.
The loss is bounded to changes made only on the offline node's disk, never propagated before it went offline.
This is verified through \texttt{TestStaleTTLReinit} in \texttt{system\_test.go}.

There are also two known limitations named sputtering cycle, determined by
Sync requests taking a CRDT state snapshot before a GC request is serviced in the CRDT actor.
This makes the Sync cycle bring along GC'd entries, creating the problem.
This is mitigated by wire-layer filtering in the \texttt{Sync} handler.
Before including a local entry in the response or applying an incoming remote entry,
the handler skips any entry where \texttt{Deleted = true} and
\texttt{now - wallTime > TTL}. Filtering is applied in both directions preventing resurrection
of tombstones regardless of GC cycles timing skew or snapshot-before-GC ordering
between nodes.

% ─────────────────────────────────────────────────────────────────────────────
\section{Lifecycle and Error Handling}
\label{sec:lifecycle}
% ─────────────────────────────────────────────────────────────────────────────

\subsection{Init / New Separation}

The CRDT package exposes three functions:
\begin{itemize}
  \item \texttt{New(workDir)} instantiates a \texttt{*CRDT} value with
    buffered channels. Pure allocation, no I/O, no error.
  \item \texttt{Init()} loads an existing node: cleans up any orphaned
    atomic-write tmp files, reads \texttt{node\_id}, and deserialises
    \texttt{crdt\_state.json}. Returns an error if the working directory
    has not been initialised (\emph{``run init first''}).
  \item \texttt{InitNew(entries)} bootstraps a fresh node: creates the
    working directory, writes \texttt{node\_id} if absent, stamps all
    supplied entries with a fresh HLC tick, and persists the resulting
    state. Used both at first-time setup and when forcing a clean slate.
\end{itemize}

The controller decides which path to take via \texttt{fixNodeState}:

\begin{lstlisting}
func (ctrl *Controller) fixNodeState(lastAlive int64) error {
    if lastAlive != -1 &&
        time.Now().UnixNano()-lastAlive > ctrl.tombstoneTTL {
        return ctrl.crdt.InitNew(types.Settings{})
    }
    return ctrl.crdt.Init()
}
\end{lstlisting}

The input is \texttt{lastAliveTime()}, defined as
$\max(last\_shutdown, last\_heartbeat)$.
\texttt{last\_shutdown} is written only by the deferred
\texttt{saveShutdownTime()} on graceful exit. \texttt{last\_heartbeat} is
(by default) written every \texttt{TombstoneTTL/4} with a dedicated goroutine in
\texttt{Run}, providing a proof-of-life signal even when the node
never reaches a graceful shutdown. Without the heartbeat the reinit
suffers two symmetric failures: a node that crashed before any
graceful shutdown would always take the \texttt{Init()} branch on restart,
even after weeks offline (cluster pollution). A node that ran for years past
its last graceful shutdown would always take the \texttt{InitNew()} branch
on a brief crash. Both files return \texttt{-1} on the first startup after
\texttt{InitNew}, going through \texttt{Init()}
on the freshly-written state.

\subsection{Error Propagation Principle}

Functions in library packages (\texttt{crdt}, \texttt{network}) return
errors rather than calling \texttt{log.Fatal} or \texttt{os.Exit} directly.
\texttt{log.Fatal} belongs only at the application boundary (\texttt{main})
where there is sufficient context to decide that an error is truly
unrecoverable. This makes every package independently testable without the test process exiting.

% ─────────────────────────────────────────────────────────────────────────────
\section{Testing Strategy}
\label{sec:testing}
% ─────────────────────────────────────────────────────────────────────────────

\subsection{Clock Injection}

The \texttt{HLC} struct accepts a \texttt{Clock} interface via parametrized
options at construction time. Tests inject \texttt{MockClock} to control
wallTime advancement deterministically, covering edge cases such as
backward clock steps, simultaneous ticks, and the case where system clock
overtakes both local and received wallTime, resetting logical counter to zero.

\subsection{Network Tests with bufconn}

The \texttt{network} package is tested with \texttt{google.golang.org/grpc/test/bufconn},
an in-memory listener that presents a real gRPC connection without opening
any network socket. This allows tests to exercise the full gRPC stack without port
allocation or network related problems.

\subsection{Coverage}

Line coverage targets were established and maintained:
\texttt{types} at 100\%, \texttt{crdt} at 90\%, \texttt{logic}
at 90\%, \texttt{network} at 91\%, and
\texttt{controller} at 91\%. The remaining uncovered branches are structurally
untestable without process-level mocking, \texttt{json.Marshal} on \texttt{map[string]string} (the error path never
fires in the Go standard library), and the \texttt{fsnotify} closed-channel
path (fires only after the watcher goroutine has already exited).
Unit tests \texttt{TestMerge\_SameWallTimeTiebreakByLogical} and
\texttt{TestMerge\_SameWallAndLogicalTiebreakByNodeID} verify the
logical-then-nodeID tiebreak defined in Section~\ref{sec:crdt}, covering
the concurrent-write scenario that system-level tests cannot exercise.

\subsection{System Tests}

System tests run against a live three-node cluster deployed with Docker Compose, each node in its own isolated container,
communicating over gRPC on the docker network. Three nodes is the minimal topology to simulate multi-peer behaviour.
A two-node setup is not enough for example to test the fan-out sync between peers.

\paragraph{Infrastructure.}
Logic helpers are built such \texttt{writeToNode} injects entries via the \texttt{Sync} RPC.
A \texttt{wallOffset} parameter shifts the
injected HLC wall time, enabling conflict-resolution scenarios
without touching the system clock.
Convergence is verified by \texttt{waitKeyValue}, \texttt{waitKeyDeleted}, and
\texttt{waitKeyPurged}, each polling all nodes until the expected condition holds
or a deadline is reached. Docker helpers
(\texttt{stopNode}, \texttt{startNode}, \texttt{waitForNodeUp}) drive
container lifecycle for fault-tolerance tests.
The compose environment overrides \texttt{tombstoneTTL} to $10s$ and
\texttt{gcInterval} to $3s$ so that GC and TTL reinit tests complete in under a
minute.

\paragraph{Basic write propagation.}
\texttt{TestInitialStateConverges} verifies that node1's seed entry
(\texttt{da=ddd}) reaches all peers through the startup anti-entropy pass.
\texttt{TestWritePropagatesAcrossNodes} injects a distinct
key on each of the three nodes and asserts cluster-wide convergence,
covering every origination point in a single test. Concurrent write tests are performed
at unit level since the actor model serializes the writes at system level.

\paragraph{Local file-watcher path.}
\texttt{TestFileWatcherPropagates} edits node1's \texttt{settings.json} from
outside the controller (via \texttt{docker exec}), bypassing the gRPC
\texttt{Sync} ingress that every other test uses. This closes the end-to-end loop on the production
write path that the unit-level \texttt{TestRun\_LocalFileChangePropagatesToCRDT}
only covers in-process.

\paragraph{LWW conflict resolution.}
\texttt{TestHLCConflictResolution} writes the same key to two nodes with
intentionally different HLC wall times (loser at $-2s$, winner at
$+500ms$) and asserts that the cluster converges to the winner.
\texttt{TestOverwriteConverges} confirms that a later write (higher HLC) on a
key that has already converged cluster-wide displaces the earlier value.
\texttt{TestOlderWriteDoesNotOverwrite} sends a stale write ($-5s$
wall offset) to a node that already holds a newer value, waits for any
incorrect propagation to surface then asserts the newer value over the cluster.

\paragraph{Tombstone lifecycle.}
\texttt{TestTombstonePropagates} writes a key and waits for convergence then deletes it via a tombstone on node1
and verifies cluster convergence.
\texttt{TestTombstoneResurrection} writes a key, deletes it, then writes the
same key again with a $+2s$ HLC offset. The cluster must accept the
resurrection since the new HLC beats the tombstone's clock. This is a guard
against tombstones permanently blocking the re-use of a key.
\texttt{TestTombstoneGC} extends \texttt{TestTombstonePropagates} by waiting
for the tombstone's wall time to age past TTL and the GC loop to fire,
confirming the entry is absent from \texttt{Sync} responses on every node.

\paragraph{Node failure recovery.}
\texttt{TestNodeRestartConverges} stops node1, writes a key on node2 while
it is offline, restarts node1 and asserts that anti-entropy delivers the
missed entry to node1.
\texttt{TestSIGKILLRecovery} proves an ungraceful crash leaves
\texttt{crdt\_state.json} intact. After a write converges, two peers are
stopped so node1 on restart cannot sync, then is killed by
\texttt{SIGKILL}. The normal refresh of \texttt{last\_shutdown} is bypassed
because of SIGKILL. The heartbeat ticker keeps \texttt{last\_heartbeat} fresh
so \texttt{lastAliveTime} on restart resolves on a recent timestamp and
\texttt{fixNodeState} takes the \texttt{Init()} branch. Node1 must still hold
the pre-crash key, proving that persistence held under a crash.
\texttt{TestStaleTTLReinit} goes through the forced-reinit path: after stopping
node1, overwrites both \texttt{last\_shutdown} and \texttt{last\_heartbeat}
with a timestamp $11s$ in the past (exceeding the $10s$ TTL), injects
a future-clock entry into node1 on-disk CRDT state (a winning entry) and
restarts. If \texttt{fixNodeState} correctly calls \texttt{InitNew}, the
injected entry is wiped before any propagation and node2's value is spread
cluster-wide, if reinit is skipped the injected entry's higher HLC beats
node2 and the test fails.

Table~\ref{tab:systemtests} summarises all system tests.

\begin{table}[htbp]
  \caption{System test suite}
  \label{tab:systemtests}
  \resizebox{\columnwidth}{!}{%
    \begin{tabular}{ll}
      \toprule
      \textbf{Test} & \textbf{Property verified} \\
      \midrule
      \texttt{TestInitialStateConverges}        & Startup anti-entropy delivers seed data \\
      \texttt{TestWritePropagatesAcrossNodes}   & Write from any node reaches all peers \\
      \texttt{TestFileWatcherPropagates}        & External \texttt{settings.json} edit propagates via fsnotify \\
      \texttt{TestHLCConflictResolution}        & Higher HLC wins conflict \\
      \texttt{TestOverwriteConverges}           & Later write displaces earlier value \\
      \texttt{TestOlderWriteDoesNotOverwrite}   & Stale write does not overwrite newer \\
      \texttt{TestTombstonePropagates}          & Delete tombstone reaches all nodes \\
      \texttt{TestTombstoneResurrection}        & Newer write defeats existing tombstone \\
      \texttt{TestTombstoneGC}                  & Tombstone purged after TTL expiry \\
      \texttt{TestNodeRestartConverges}         & Offline node catches missed updates \\
      \texttt{TestSIGKILLRecovery}              & State survives ungraceful crash \\
      \texttt{TestStaleTTLReinit}               & Long-offline node forces clean reinit \\
      \bottomrule
    \end{tabular}%
  }
\end{table}

% ─────────────────────────────────────────────────────────────────────────────
\section{Conclusion}
\label{sec:conclusion}
% ─────────────────────────────────────────────────────────────────────────────

This paper presented a production-quality CRDT-based configuration
synchronisation system. The core LWW CRDT provides convergence
under concurrent writes and network partitions. The actor concurrency model
eliminates lock contention. Anti-entropy reconciliation bounds
divergence to Sync interval regardless of broadcast reliability.
Atomic persistence through \texttt{os.Rename} prevents corruption on ungraceful shutdown.
Tombstone GC fixes CRDT state growth for long-running deployments.

Several design decisions evolved during implementation. The protobuf
abstraction boundary was moved into the \texttt{network} package, removing
transport-layer knowledge from the controller. Error propagation got modified
to return \texttt{error} from library code and keep
\texttt{log.Fatal} in the application entry-point. Forceful init was added
for nodes rejoining the cluster after more than TTL elapsed since last shutdown. A full
reconciliation cycle was added upon start to make sure ungraceful shutdowns did not
produce stale entries.

The system does not address gossip-based broadcast efficiency, or metrics
observability, gaps appropriate for future work in production
deployments at larger scale.

% ─────────────────────────────────────────────────────────────────────────────
\section*{Acknowledgments}
% ─────────────────────────────────────────────────────────────────────────────

This project used a coding assistant during development and its use is disclosed here.
The final implementation, architectural choices and system design were all.
The assistant role was to review, be a sounding board and a template creator for the code and report.
All its output was verified and edited before being incorporated.

Specifically it was used to:
\begin{itemize}
  \item \textbf{Clarify theory.} Discuss and pressure-test the author's understanding
    of the CRDT formalism used in this report: join-semilattices, the LWW-register
    model, least-upper-bound merge semantics and proof correctness so that the claims
    made here are stated precisely.
  \item \textbf{Review the manuscript.} Identify imprecise or overstated claims,
    grammatical errors, and notation issues in the prose, suggest phrasing.
    Additions were made by the assistant then edited and reviewed by the author.
  \item \textbf{Review code and find defects.} Audit selected portions of the code.
    This surfaced a write-ordering bug that could have caused out of order writes.
  \item \textbf{Format references.} Produce and check bibliographic entries for
    citations in the paper.
  \item \textbf{Boilerplate function creation.} Many unit tests were coded by the assistant to save time.
    All generated tests were reviewed in a later refactor. Redundant or incorrect tests were removed.
    System tests and specific test scenarios were fully planned and made by hand.
    Channel accessor boilerplates were templated by the assistant (Get, Snapshot, NotifyRemote, etc.).
\end{itemize}

Assistant's role was limited to review, suggestions, and the specific contributions listed above.

% ─────────────────────────────────────────────────────────────────────────────
\begin{thebibliography}{00}

  \bibitem{raft}
  D.~Ongaro and J.~Ousterhout,
  ``In search of an understandable consensus algorithm,''
  in \emph{Proc. USENIX Annual Technical Conf. (USENIX ATC)},
  Philadelphia, PA, Jun.\ 2014, pp.\ 305--319.

  \bibitem{gilbert2012}
  S.~Gilbert and N.~A.~Lynch,
  ``Perspectives on the CAP theorem,''
  \emph{Computer}, vol.~45, no.~2, pp.~30--36, Feb.\ 2012.

  \bibitem{shapiro2011}
  M.~Shapiro, N.~Pregui\c{c}a, C.~Baquero, and M.~Zawirski,
  ``A comprehensive study of convergent and commutative replicated data types,''
  INRIA, Tech. Rep. RR-7506, 2011.

  \bibitem{lamport1978}
  L.~Lamport,
  ``Time, clocks, and the ordering of events in a distributed system,''
  \emph{Commun. ACM}, vol.~21, no.~7, pp.~558--565, Jul.\ 1978.

  \bibitem{kulkarni2014}
  S.~S.~Kulkarni and M.~Demirbas,
  ``Logical physical clocks and consistent snapshots in globally distributed
  databases,''
  in \emph{Proc. 18th Int. Conf. Principles of Distributed Systems (OPODIS)},
  Cortina d'Ampezzo, Italy, Dec.\ 2014.

  \bibitem{cockroachdb2023}
  CockroachDB Engineering,
  ``Living without atomic clocks,''
  Cockroach Labs Blog, 2023.
  [Online]. Available: \url{https://www.cockroachlabs.com/blog/living-without-atomic-clocks/}

  \bibitem{dynamo2007}
  G.~DeCandia, D.~Hastorun, M.~Jampani, G.~Kakulapati, A.~Lakshman,
  A.~Pilchin, S.~Sivasubramanian, P.~Vosshall, and W.~Vogels,
  ``Dynamo: Amazon's highly available key-value store,''
  in \emph{Proc. 21st ACM Symp. Operating Systems Principles (SOSP)},
  Stevenson, WA, Oct.\ 2007, pp.\ 205--220.

\end{thebibliography}

\end{document}
